%! Parent Functions and Transformations — Unit 2, Lesson 2.2 — Warm-Up
\documentclass[10pt]{article}
\usepackage{atda-article}
\usepackage{atda-boxes}
\newcommand{\TallMath}[1]{$\displaystyle #1\rule[-1.4em]{0pt}{3.2em}$}
\setlength{\workrowsep}{6pt}   % handwriting room; moves blank and key together

\begin{document}

\pageheader{Unit 2, Lesson 2.2}{Warm-Up}
% NAMESTRIP: no \namedateperiod here. The student writes their name once, on the
% cover this component is stapled behind. See references/conventions.md.

\vspace{0.15in}

\begin{notesbox}{Where Does the Number Go In?}
\small
Five quick items. Items 2 and 3 look almost the same and are not --- read them slowly.

\begin{enumerate}[leftmargin=1.6em, itemsep=4pt, topsep=4pt]

\item \textbf{Two rules, one table.} Fill in every cell.
\smallskip

\renewcommand{\arraystretch}{1.3}
\begin{tabular}{@{}|c|c|c|c|c|@{}}
\hline
$x$ & $0$ & $1$ & $2$ & $3$ \\ \hline
$f(x)=x^2$ & \blank{1.1cm} & \blank{1.1cm} & \blank{1.1cm} & \blank{1.1cm} \\ \hline
$h(x)=2^x$ & \blank{1.1cm} & \blank{1.1cm} & \blank{1.1cm} & \blank{1.1cm} \\ \hline
\end{tabular}
\smallskip

Now find $h(-1)$ and $h(-2)$ (Lesson 1.1 --- a negative exponent).
\begin{work}
  h(-1) &= 2^{-1} \\
        &= \frac{1}{2} \\
  h(-2) &= 2^{-2} \\
        &= \frac{1}{4}
\end{work}

\item \textbf{Feeding the rule a different number.} Let $f(x)=x^2$. Find $f(5)$, then find
$f(5-3)$.
\begin{work}
  f(5) &= 5^2 \\
       &= 25 \\
  f(5-3) &= f(2) \\
         &= 4
\end{work}

\item \textbf{Where the minus sign sits.} Evaluate both of these at $x=3$.
\begin{work}
  -x^2 &= -(3)^2 \\
       &= -9 \\
  (-x)^2 &= (-3)^2 \\
         &= 9
\end{work}
Are $-x^2$ and $(-x)^2$ the same rule? Say how you know.

\writelines{1}

\item \textbf{Writing a set down} (Lesson 2.1). \quad (a) Write $y \ge 3$ in interval notation.
\quad (b) Write $(-\infty, 0]$ as an inequality.

\writelines{1}

\item \textbf{Last night's homework, in one sentence.} You found that every output of
$g(x)=x^2+3$ is exactly $3$ more than the matching output of $f(x)=x^2$. Without drawing
anything, say how the graph of $g$ sits compared to the graph of $f$.

\writelines{1}

\end{enumerate}
\end{notesbox}

\end{document}
